\section{Introduction}
\label{sec:introduction}

The simplicial volume of a closed connected oriented \(n\)-manifold
\(M\) is the \(\ell^1\)-seminorm of its real fundamental class:
\[
\|M\|=
\inf\left\{\sum_j|a_j|:
 \sum_j a_j\sigma_j\text{ is a finite real singular cycle representing }
 [M]_{\R}\right\}.
\]
It measures topological complexity through the cost of a fundamental
cycle. Gromov's work on bounded cohomology connected this invariant
with Riemannian volume under lower Ricci curvature bounds
\cite{Gromov1982}. A lower scalar curvature bound retains much less
directional information. The possibility that it still controls
simplicial volume is therefore a particularly strong relation between
local geometry and global topology.

In 1986, Gromov proposed precisely the integral estimate considered
here: Conjecture 3.A, Remark (B), Equation (8\('\)) of
\cite{Gromov1986} asks for a dimension-dependent upper bound on
simplicial volume by the integral of the \(n/2\)-power of the negative
part of scalar curvature. With
\[
  \Scal_g^-:=\max\{0,-\Scal_g\},
\]
the integral
\(\int_M(\Scal_g^-)^{n/2}\,dV_g\) is invariant under constant rescaling
of \(g\). It records both the size and the concentration of negative
scalar curvature.

The conjecture is closely related to the uniform volume assertion
\[
\Scal_g\ge-\kappa^2
\quad\Longrightarrow\quad
\|M\|\le C_n\kappa^n\Vol_g(M),
\qquad \kappa\ge0.
\]
Although the integral formulation keeps track of a varying scalar
curvature, the two assertions, when required for every manifold and
metric in a fixed dimension \(n\ge3\), are equivalent by the Yamabe
reformulation. A precise statement appears in
\cite[Lemma 2.1]{MinZhengZhu2026}.

Several results establish related forms under additional geometric or
topological hypotheses. Braun and Sauer prove a macroscopic counterpart
when the volumes of unit balls in the universal cover have a uniform
upper bound; their constant also depends on that bound
\cite[Theorem 1.1]{BraunSauer2021}. For the qualitative case of
nonnegative scalar curvature, Ma, Wang, Xie, Yu, and Zhu prove
vanishing when the universal cover is spin
\cite{MaWangXieYuZhu2026}. Min, Zheng, and Zhu prove the integral
estimate for every Riemannian metric on every closed K\"ahler surface,
with coefficient \(27/2\) in the upper bound for \(\|M\|\), and also
treat an infinite family of non-K\"ahler symplectic four-manifolds
\cite[Theorem 1.2 and Section 4]{MinZhengZhu2026}.

Our argument uses the vanishing theorem of the companion manuscript
\cite[Corollary 1.5]{OpenAIScalarInessentiality2026}: a closed connected
oriented smooth manifold of positive dimension admitting a metric of
nonnegative scalar curvature has zero
real simplicial volume. We state the exact input in
Theorem~\ref{thm:vanishing-input}. The work of this paper is the passage from
that qualitative vanishing result to a uniform quantitative estimate.

\begin{theorem}
\label{thm:main}
For every integer \(n\ge3\), there is a constant \(a_n>0\), depending
only on \(n\), such that every closed connected oriented smooth
\(n\)-manifold \(M\) and every smooth Riemannian metric \(g\) on \(M\)
satisfy
\begin{equation}
  \int_M(\Scal_g^-)^{n/2}\,dV_g\ge a_n\|M\|.
  \label{eq:main}
\end{equation}
\end{theorem}

For example, a metric with \(\Scal_g\ge-n(n-1)\) satisfies
\[
  \Vol_g(M)\ge
  a_n[n(n-1)]^{-n/2}\|M\|.
\]
The constant is uniform over the topology of \(M\) and over all its
metrics. In particular, \eqref{eq:main} permits the negative scalar
curvature to be concentrated in a small region.

\subsection*{The route to the estimate}

The first step is a conformal reduction. If \(\|M\|>0\), the vanishing
input forces the Yamabe minimizer in every conformal class to have
negative scalar curvature. Normalize it to scalar curvature \(-1\).
H\"older's inequality bounds the volume of this normalized metric by
the integral in \eqref{eq:main} for the original metric. It remains
to prove
\[
  \|M\|\le C_n\Vol_g(M)\qquad\text{when }\Scal_g=-1.
\]
This reduction is given in Proposition~\ref{prop:yamabe-reduction}. The
quantitative argument begins with this normalized problem.

Bounded cohomology supplies a cocycle of supremum norm at most one
whose value \(I\) on the fundamental class is at least
\(\|M\|/2\). We realize it as a closed differential form on a simplex
of probability vectors. An equivariant probability map from the
universal cover pulls this form back with integral \(I\). Thus the
remaining task is to control this integral by the original volume.

The first construction changes the metric and the probability map
together. For a metric \(h\) and weight \(\phi\), the relevant scalar
quantity is
\[
  \wscal(h,\phi)=\Scal_h-2\Delta_h\phi-|d\phi|_h^2.
\]
With \(\phi=-f\), this is Perelman's modified scalar curvature for the
measure \(e^{-f}dV_h\) \cite[Section 1.3]{Perelman2002}.
A graph deformation enlarges \(h\) while changing \(\phi\) so that
\(e^\phi dV_h\) is preserved. Its scalar identity contains positive
squares and a derivative of the prescribed weighted mean curvature.
This adapts the graph stabilization and prescribed-curvature method of
\cite[Proposition 4.1 and Theorem 5.2]{OpenAIScalarInessentiality2026}.
Choosing that prescription in terms of one probability coordinate at
a time gives a gain in the energy of the probability map. Iterating
the construction produces a measured equivariant family over \(M\),
with a density \(\rho\le dV_g\), a new metric and weight, and a large
coefficient \(E\) satisfying the leafwise inequality
\[
  \wscal(h,\phi)\ge
  -C_n+2\log\frac{dV_h}{\rho}-n\log E+E|dF|_{h,\HS}^2.
\]
Integrals over this family include its auxiliary probability measure.
The logarithmic volume term and the energy term have distinct roles:
the first will cancel the volume expansion, and the second will pay
for derivatives of comparison coordinates. The construction also
makes \(|dF|_{h,\HS}\) small in absolute terms. This allows its label
overlaps to be controlled through a fixed graph even after many
deformations. Theorem~\ref{thm:entropy} states the full result, including the
relation between \(E\) and the complexity of that graph.

The second construction estimates the degree of a map from a
measured family of \(d\)-manifolds to a \(d\)-cube. Its bound contains
\(-\wscal/2\) and a separate penalty for the derivative of each
coordinate. The perimeter-minus-forcing functional is the
\(\mu\)-bubble functional used in scalar curvature geometry
\cite{Gromov1996,ChodoshLi2024}. Here an ordered family of its
minimizing boundaries reduces the dimension one coordinate at a time.
Their normal speed, or lapse, appears both in the scalar
identity and in coarea, producing the needed cancellation. In high
dimensions these boundaries can be singular. We move the residual
map away from the target near the singular region, while weighted
distance potentials supply the scalar gain needed to control the
change. The resulting Theorem~\ref{thm:degree} is formulated for general
controlled measured bands; it does not depend on simplicial volume.

The final construction connects the cocycle integral \(I\) to these
degrees. Integrating the cocycle form on small primal grid cells gives
a balanced chain of complementary dual cubes. This follows the
probability-map and auxiliary-degree strategy of
\cite[Section 3]{OpenAIScalarInessentiality2026}. An unbiased rounding
preserves the balance relations with integer multiplicities. We adapt
Gaifullin's reflected permutohedral realization
\cite{Gaifullin2008,Gaifullin2013}, equipping the pairing choices with
normalized probability laws and controlling the resulting smooth
charts. An oppositely oriented mirror copy closes the unmatched faces.
In the neighborhood used for comparison, the resolved family retains a
side whose oriented pushforward is exactly the original rounded chain.
The construction supplies the chart and metric bounds needed for the
degree estimate and the final averaging.

Compare the probability map with the resolved dual cubes in dimension
\(q=n+m\), where \(m\) is the dimension of a dual cube. Let
\(h,\rho,E\) be the deformed metric, retained density, and energy
coefficient produced above, and set \(W=\log(dV_h/\rho)\). Assign the
comparison lengths \(L_T=b_*\sqrt E\) in the \(m\) directions tangent
to the cube and \(L_N=b_*\sqrt{qE}\) in the \(n\) normal directions,
with \(b_*>0\) fixed dimensionally. The degree estimate contributes
\((L_T^mL_N^n)^{-1}\), the product metric contributes \(L_T^m\), and
the scalar term \(2W-n\log E\) contributes \(E^{n/2}e^{-W}\).
The factors involving \(E\) and the volume density therefore satisfy
\[
 \frac{L_T^m}{L_T^mL_N^n}E^{n/2}e^{-W}dV_h
   =b_*^{-n}q^{-n/2}\rho.
\]
On the straight part, where the resolution metric is pulled back from the
flat cube, averaging the normal derivative penalties together with the
Hilbert bound for the cocycle form costs at most \(C_nq^{n/2}\). For this
contribution, if \(r\) is the grid spacing, the cell coefficient,
dual-cube volume, and density of grid centers contribute
\(r^n r^m r^{-q}=1\). The seam contribution is controlled by its small
coordinate volume. The required lower bound for \(E\) is an elementary
function of \(q\), bounded by an exponential tower whose fixed height and
base polynomial depend only on \(n\). It is fixed by this transfer construction before
the entropy theorem is applied. These balances bound the detector pairing by
\(I\le C_n\int\rho\le C_n\Vol_g(M)\).

Section~\ref{sec:reduction-detector} gives the conformal reduction and the
detector. Section~\ref{sec:graphs} establishes the graph identity, the
prescribed graph equations, and the estimate for one sweep of the
probability coordinates. Section~\ref{sec:entropy} iterates that estimate and
controls the label graph. Section~\ref{sec:degree} proves the degree theorem,
including the treatment of singular hypersurfaces.
Section~\ref{sec:transfer} constructs the dual chains and their smooth
resolution and completes the proof of Theorem~\ref{thm:main}.
